\documentclass[reqno]{amsart}
\usepackage[T1]{fontenc}
\usepackage{amsmath,amssymb,amsthm}
\usepackage{xr-hyper}
\usepackage[hidelinks]{hyperref}
\externaldocument[da-]{20261126T030300--the-data__workbook}
\externaldocument[bp-]{20261126T020210--bessel-poincare-series-and-their-gram-operators__workbook}
\externaldocument[m1-]{20261126T051900--unfolding-twice-gram-operators-as-sums-over-rational-points__workbook}
\externaldocument[m2-]{20261126T052000--which-terms-survive-characters-minors-and-the-gain-at-the-level__workbook}

\newcounter{chap}
% Chapter numbers, generated from 20261125T012000--grh-workbook__guide.org
% by scripts/make-chapters.py.  Do not edit: change the order of the
% index and rerun the script (./preview-site.sh and CI do this).
%
% Usage in a sheet, after \newcounter{chap}:
%   % Chapter numbers, generated from 20261125T012000--grh-workbook__guide.org
% by scripts/make-chapters.py.  Do not edit: change the order of the
% index and rerun the script (./preview-site.sh and CI do this).
%
% Usage in a sheet, after \newcounter{chap}:
%   % Chapter numbers, generated from 20261125T012000--grh-workbook__guide.org
% by scripts/make-chapters.py.  Do not edit: change the order of the
% index and rerun the script (./preview-site.sh and CI do this).
%
% Usage in a sheet, after \newcounter{chap}:
%   \input{chapters}\chapnumber{<sheet ID>}
\makeatletter
\providecommand\chapters@set[2]{\@namedef{chapters@#1}{#2}}
\chapters@set{20261125T001000}{1}% zeros-to-ramp-sums
\chapters@set{20261125T001100}{2}% inert-primes-and-estimate-b
\chapters@set{20261125T005000}{3}% galois-fourier-and-averaging
\chapters@set{20261125T005500}{4}% columns-from-a-sphere
\chapters@set{20261125T000000}{5}% slope-method
\chapters@set{20261126T030100}{6}% a-twist-for-each-prime
\chapters@set{20261125T003000}{7}% the-chain-and-the-square-root
\chapters@set{20261125T010000}{8}% rank-lemma-and-zero-estimate
\chapters@set{20261125T010500}{9}% poincare-series-and-near-orthogonality
\chapters@set{20261124T223000}{10}% elliptic-curves-for-the-proof
\chapters@set{20261125T009500}{11}% fock-model-and-theta-intertwiners
\chapters@set{20261124T223100}{12}% complex-multiplication-and-inert-reduction
\chapters@set{20261125T008000}{13}% one-cm-determinant
\chapters@set{20261124T214500}{14}% supersingular-theta-divisibility
\chapters@set{20261124T214600}{15}% the-isotropic-plane
\chapters@set{20261125T013700}{16}% pairings-on-a-formal-group-of-height-two
\chapters@set{20261125T013500}{17}% hecke-nondegeneracy-on-the-tree
\chapters@set{20261126T030200}{18}% the-bound-is-attained
\chapters@set{20261125T009000}{19}% kronecker-limit-and-chowla-selberg
\chapters@set{20261125T002000}{20}% definite-theta-lifts
\chapters@set{20261125T015500}{21}% siegel-weil-in-the-smallest-case
\chapters@set{20261125T002100}{22}% definite-theta-lifts-and-norms
\chapters@set{20261125T016500}{23}% hecke-operators-and-theta-lifts
\chapters@set{20261125T016700}{24}% divisibility-along-the-chain
\chapters@set{20261125T004000}{25}% denominators-and-sha
\chapters@set{20261125T004100}{26}% tate-twists-and-sha
\chapters@set{20261126T020430}{27}% the-minus-part-by-gauss-sums
\chapters@set{20261125T014100}{28}% congruence-modules
\chapters@set{20261125T014500}{29}% eisenstein-congruences-and-selmer-classes
\chapters@set{20261125T014700}{30}% pseudocharacters-and-the-length-bound
\chapters@set{20261125T011000}{31}% galois-representations-black-box
\chapters@set{20261126T030300}{32}% the-data
\chapters@set{20261126T020210}{33}% bessel-poincare-series-and-their-gram-operators
\chapters@set{20261126T030400}{34}% the-seesaw-identity
\chapters@set{20261126T030500}{35}% divisibility-of-the-composite-lift
\chapters@set{20261126T030600}{36}% theorem-a-short-forms
\chapters@set{20261126T030700}{37}% the-determinant-of-a-chain
\chapters@set{20261126T030800}{38}% the-cost-of-one-lift
\chapters@set{20261126T030900}{39}% denominators-of-projections-to-lifts
\chapters@set{20261126T031000}{40}% twists-with-few-congruences
\chapters@set{20261126T031100}{41}% theorem-c-and-the-main-inequality
\providecommand\chapnumber[1]{%
  \@ifundefined{chapters@#1}%
    {\PackageError{chapters}{Sheet #1 is not listed in the index}%
      {List it in 20261125T012000--grh-workbook__guide.org\MessageBreak
       and run scripts/make-chapters.py.}}%
    {\edef\thechap{\@nameuse{chapters@#1}}}}
\makeatother
\chapnumber{<sheet ID>}
\makeatletter
\providecommand\chapters@set[2]{\@namedef{chapters@#1}{#2}}
\chapters@set{20261125T001000}{1}% zeros-to-ramp-sums
\chapters@set{20261125T001100}{2}% inert-primes-and-estimate-b
\chapters@set{20261125T005000}{3}% galois-fourier-and-averaging
\chapters@set{20261125T005500}{4}% columns-from-a-sphere
\chapters@set{20261125T000000}{5}% slope-method
\chapters@set{20261126T030100}{6}% a-twist-for-each-prime
\chapters@set{20261125T003000}{7}% the-chain-and-the-square-root
\chapters@set{20261125T010000}{8}% rank-lemma-and-zero-estimate
\chapters@set{20261125T010500}{9}% poincare-series-and-near-orthogonality
\chapters@set{20261124T223000}{10}% elliptic-curves-for-the-proof
\chapters@set{20261125T009500}{11}% fock-model-and-theta-intertwiners
\chapters@set{20261124T223100}{12}% complex-multiplication-and-inert-reduction
\chapters@set{20261125T008000}{13}% one-cm-determinant
\chapters@set{20261124T214500}{14}% supersingular-theta-divisibility
\chapters@set{20261124T214600}{15}% the-isotropic-plane
\chapters@set{20261125T013700}{16}% pairings-on-a-formal-group-of-height-two
\chapters@set{20261125T013500}{17}% hecke-nondegeneracy-on-the-tree
\chapters@set{20261126T030200}{18}% the-bound-is-attained
\chapters@set{20261125T009000}{19}% kronecker-limit-and-chowla-selberg
\chapters@set{20261125T002000}{20}% definite-theta-lifts
\chapters@set{20261125T015500}{21}% siegel-weil-in-the-smallest-case
\chapters@set{20261125T002100}{22}% definite-theta-lifts-and-norms
\chapters@set{20261125T016500}{23}% hecke-operators-and-theta-lifts
\chapters@set{20261125T016700}{24}% divisibility-along-the-chain
\chapters@set{20261125T004000}{25}% denominators-and-sha
\chapters@set{20261125T004100}{26}% tate-twists-and-sha
\chapters@set{20261126T020430}{27}% the-minus-part-by-gauss-sums
\chapters@set{20261125T014100}{28}% congruence-modules
\chapters@set{20261125T014500}{29}% eisenstein-congruences-and-selmer-classes
\chapters@set{20261125T014700}{30}% pseudocharacters-and-the-length-bound
\chapters@set{20261125T011000}{31}% galois-representations-black-box
\chapters@set{20261126T030300}{32}% the-data
\chapters@set{20261126T020210}{33}% bessel-poincare-series-and-their-gram-operators
\chapters@set{20261126T030400}{34}% the-seesaw-identity
\chapters@set{20261126T030500}{35}% divisibility-of-the-composite-lift
\chapters@set{20261126T030600}{36}% theorem-a-short-forms
\chapters@set{20261126T030700}{37}% the-determinant-of-a-chain
\chapters@set{20261126T030800}{38}% the-cost-of-one-lift
\chapters@set{20261126T030900}{39}% denominators-of-projections-to-lifts
\chapters@set{20261126T031000}{40}% twists-with-few-congruences
\chapters@set{20261126T031100}{41}% theorem-c-and-the-main-inequality
\providecommand\chapnumber[1]{%
  \@ifundefined{chapters@#1}%
    {\PackageError{chapters}{Sheet #1 is not listed in the index}%
      {List it in 20261125T012000--grh-workbook__guide.org\MessageBreak
       and run scripts/make-chapters.py.}}%
    {\edef\thechap{\@nameuse{chapters@#1}}}}
\makeatother
\chapnumber{<sheet ID>}
\makeatletter
\providecommand\chapters@set[2]{\@namedef{chapters@#1}{#2}}
\chapters@set{20261125T001000}{1}% zeros-to-ramp-sums
\chapters@set{20261125T001100}{2}% inert-primes-and-estimate-b
\chapters@set{20261125T005000}{3}% galois-fourier-and-averaging
\chapters@set{20261125T005500}{4}% columns-from-a-sphere
\chapters@set{20261125T000000}{5}% slope-method
\chapters@set{20261126T030100}{6}% a-twist-for-each-prime
\chapters@set{20261125T003000}{7}% the-chain-and-the-square-root
\chapters@set{20261125T010000}{8}% rank-lemma-and-zero-estimate
\chapters@set{20261125T010500}{9}% poincare-series-and-near-orthogonality
\chapters@set{20261124T223000}{10}% elliptic-curves-for-the-proof
\chapters@set{20261125T009500}{11}% fock-model-and-theta-intertwiners
\chapters@set{20261124T223100}{12}% complex-multiplication-and-inert-reduction
\chapters@set{20261125T008000}{13}% one-cm-determinant
\chapters@set{20261124T214500}{14}% supersingular-theta-divisibility
\chapters@set{20261124T214600}{15}% the-isotropic-plane
\chapters@set{20261125T013700}{16}% pairings-on-a-formal-group-of-height-two
\chapters@set{20261125T013500}{17}% hecke-nondegeneracy-on-the-tree
\chapters@set{20261126T030200}{18}% the-bound-is-attained
\chapters@set{20261125T009000}{19}% kronecker-limit-and-chowla-selberg
\chapters@set{20261125T002000}{20}% definite-theta-lifts
\chapters@set{20261125T015500}{21}% siegel-weil-in-the-smallest-case
\chapters@set{20261125T002100}{22}% definite-theta-lifts-and-norms
\chapters@set{20261125T016500}{23}% hecke-operators-and-theta-lifts
\chapters@set{20261125T016700}{24}% divisibility-along-the-chain
\chapters@set{20261125T004000}{25}% denominators-and-sha
\chapters@set{20261125T004100}{26}% tate-twists-and-sha
\chapters@set{20261126T020430}{27}% the-minus-part-by-gauss-sums
\chapters@set{20261125T014100}{28}% congruence-modules
\chapters@set{20261125T014500}{29}% eisenstein-congruences-and-selmer-classes
\chapters@set{20261125T014700}{30}% pseudocharacters-and-the-length-bound
\chapters@set{20261125T011000}{31}% galois-representations-black-box
\chapters@set{20261126T030300}{32}% the-data
\chapters@set{20261126T020210}{33}% bessel-poincare-series-and-their-gram-operators
\chapters@set{20261126T030400}{34}% the-seesaw-identity
\chapters@set{20261126T030500}{35}% divisibility-of-the-composite-lift
\chapters@set{20261126T030600}{36}% theorem-a-short-forms
\chapters@set{20261126T030700}{37}% the-determinant-of-a-chain
\chapters@set{20261126T030800}{38}% the-cost-of-one-lift
\chapters@set{20261126T030900}{39}% denominators-of-projections-to-lifts
\chapters@set{20261126T031000}{40}% twists-with-few-congruences
\chapters@set{20261126T031100}{41}% theorem-c-and-the-main-inequality
\providecommand\chapnumber[1]{%
  \@ifundefined{chapters@#1}%
    {\PackageError{chapters}{Sheet #1 is not listed in the index}%
      {List it in 20261125T012000--grh-workbook__guide.org\MessageBreak
       and run scripts/make-chapters.py.}}%
    {\edef\thechap{\@nameuse{chapters@#1}}}}
\makeatother

\ifcsname chapters@20261126T080300\endcsname
  \chapnumber{20261126T080300}
\else
  \renewcommand{\thechap}{?}
\fi
\renewcommand{\thesection}{\thechap.\arabic{section}}
\theoremstyle{definition}
\newtheorem{problem}{Problem}[section]
\newtheorem{definition}{Definition}[section]
\theoremstyle{plain}
\newtheorem{theorem}{Theorem}[section]
\theoremstyle{remark}
\newtheorem*{leanremark}{In the Lean formalization}
\newcommand{\Z}{{\mathbb{Z}}}
\newcommand{\C}{{\mathbb{C}}}
\newcommand{\ii}{\mathrm{i}}
\newcommand{\Oc}{\mathcal O}
\newcommand{\Uc}{{\mathcal{U}}}
\newcommand{\Bc}{\mathcal B}
\newcommand{\Rb}{\bar R_0}
\newcommand{\Sf}{\mathfrak S}
\newcommand{\Amp}{\mathrm{Amp}}
\newcommand{\Did}{D_{\mathrm{id}}}
\newcommand{\Dmain}{D_{\mathrm{main}}}
\newcommand{\Dcell}{D_{\mathrm{cell}}}
\newcommand{\disc}{\operatorname{disc}}
\newcommand{\lk}{\operatorname{lk}}
\newcommand{\mt}{\operatorname{mt}}
\newcommand{\ts}{\mathrm{ts}}
\newcommand{\pc}{\mathrm{pc}}
\newcommand{\sz}{\operatorname{sz}}
\newcommand{\calb}{\operatorname{cal}}
\newcommand{\cri}{\operatorname{cr}}
\newcommand{\dsc}{\operatorname{dsc}}
\newcommand{\crk}{\operatorname{crk}}
\newcommand{\lcr}{\operatorname{lcr}}
\newcommand{\Rep}{\mathrm{Rep}}
\newcommand{\Nm}{\operatorname{Nm}}
\newcommand{\diag}{\operatorname{diag}}
\newcommand{\A}{\mathbb A}
\newcommand{\tl}[1]{\texttt{#1}}
\setlength{\emergencystretch}{3em}

\title{Size, gain and count: the remaining terms}
\author{}
\date{November 26, 2026}

\begin{document}
\maketitle
\refstepcounter{chap}\label{chap}

\paragraph{Prerequisites.}
Chapter~\ref{m2-chap} in full, in particular the frame \(v_1,\dots,v_\varrho,v_{\varrho+1}:=z\), the leading minors \(\langle\gamma\rangle_j\), \(\disc\), \(\lk\), the remaining terms \(\Gamma^*\), \(\mt\), the second moments and the link second moment (Definitions~\ref{m2-def:frame}--\ref{m2-def:meet} and~\ref{m2-def:moments}), Theorems~\ref{m2-thm:survivors}, \ref{m2-thm:link-bound}, \ref{m2-thm:cut} and Problems~\ref{m2-pr:cross}, \ref{m2-pr:columns}, \ref{m2-pr:hua}.  From Chapter~\ref{m1-chap}: amplitude sums (Definition~\ref{m1-def:amplitude}; \(\Amp^\flat\) is formed with the cut function), \(\Did\), \(\Dmain\) and, for the data, \(b^{11kf}\Did\le\Dmain\le(M/b^6)^{kf}\Did\).  From Chapter~\ref{bp-chap}: classes \(c\) and representatives \(h_c\), integrality and reduction, \(K_M\), \(\Uc\), \(\nu\) (Definition~\ref{bp-def:scaled}), \(\Rb\) (Definition~\ref{bp-def:bessel-rep}), admissible fillings (Definition~\ref{bp-def:filling}).  From Definition~\ref{da-def:data}: \(b\ge1000\), \(b^{10}\le p_i\le b^{11}\), \(M=p_0p_1\), and \(\varrho+1=m/2<b\) (Problem~\ref{da-pr:chain} and the definition of \(\varrho\) after Problem~\ref{da-pr:signs}); \(kf\ge40\) (Chapter~\ref{m2-chap}).  Hermitian forms are linear in the first variable; \(\lVert\cdot\rVert\) is the Euclidean norm of the entries; \(k'\ge1\) is an integer (primed, since \(k\) belongs to the data).

\medskip
This sheet serves the second moment in the proof of Theorem~\ref{bp-thm:B} of Chapter~\ref{bp-chap}, \emph{Theorem~B: enough forms}.  Chapter~\ref{m2-chap} states two bounds.  Theorem~\ref{m2-thm:link-bound} bounds the sum of the squared norms of the cut blocks of the terms of nonzero discriminant, taken separately for each link, each pair of classes and each reduction in the specified set.  Theorem~\ref{m2-thm:cut} bounds the square root of the full second moment of the terms off the Bessel subgroup over sets that avoid the scalar cosets.  There the amplitude sums in each block agree with those of \(\Gamma^*\) by Theorem~\ref{m2-thm:survivors}.  What are they made of?  The norm of a cut block of a stable set of remaining terms is at most \(\nu\) times the sum of the \emph{sizes} of its near representatives.  For two terms with cell ranks over one block with agreeing reduction, the \emph{gain} bounds the discriminant of their common lattice by their \emph{calibres}, which enter the sizes inversely.  For pairs of near representatives of fixed corank and bounded discriminant, the \emph{count} bounds the sum of the inverse \(k\)-th powers of the discriminants.  For the data, we derive both bounds of Chapter~\ref{m2-chap} using the size, gain, count and tail estimates of \cite{GL}, together with the stability and summation properties of the pieces in \S\ref{sec:back}.

\section{The size at one place}\label{sec:one-place}

The next four problems compute, at one place, a weighted Cauchy--Schwarz inequality, the mean of a lattice sum, and a determinant integral.

\begin{problem}\label{pr:schur}
For finitely many \(x_j,y_j\ge0\) and \(v_j>0\), \((\sum_jx_jy_j)^2\le(\sum_jx_j^2/v_j)(\sum_jy_j^2v_j)\).  For \(x=(1,2)\), \(y=(3,1)\), \(v=(2,\tfrac12)\): \(25\le\tfrac{17}2\cdot\tfrac{37}2=\tfrac{629}4\).
\end{problem}

\begin{problem}\label{pr:mean}
Let \(L\subset\C^{k'}\) be a lattice of rank \(2k'\) and \(\mathsf h\ge0\) integrable.  The mean, for the Haar probability measure on \(\C^{k'}/L\), of \(\sum_{v\in L}\mathsf h(a+v)\) is \(\int\mathsf h/\operatorname{covol}(L)\).  For \(\mathsf h(y)=(1+\lvert y\rvert^2)^{-k'-1}\) it is \(\pi^{k'}/(k'!\operatorname{covol}(L))\); for \(L=(2+\ii)\Z[\ii]^2\) it is \(\pi^2/50\approx0.197\).
\end{problem}

\begin{problem}\label{pr:covolume}
A lattice of rank \(n'\) in \(\mathbb R^{n'}\) whose inner products are integers has covolume at least \(1\).
\end{problem}

\begin{problem}\label{pr:hua-row}
Let \(\alpha\in\C^\times\), \(y_0\in\C^{k'}\) a row, \(A\) a positive definite hermitian matrix of size \(k'\), and \(v:=\alpha y+y_0\) for rows \(y\in\C^{k'}\).  Then
\[
 \int_{\C^{k'}}\det(A+v^*v)^{-k'-1}\,dy=\lvert\alpha\rvert^{-2k'}(\det A)^{-k'}\,\frac{\pi^{k'}}{k'!}.
\]
For \(k'=7\), \(\alpha=2\), \(A=\diag(3,1,\dots,1)\) the value is \(\pi^7/(7!\,2^{14}3^7)\approx1.672\cdot10^{-8}\).
\end{problem}

For \(i'\) rows Hua's bound holds: if \(k'\ge7\), \(\alpha\) is an invertible complex \(i'\times i'\) matrix, \(y_0\) a complex \(i'\times k'\) matrix and \(A\) positive definite of size \(k'\), then \(\int\lvert\det(A+v^*v)\rvert^{-k'-i'}\,dy\le\lvert\det\alpha\rvert^{-2k'}(\det A)^{-k'}\) with \(v:=\alpha y+y_0\), the integral over the complex \(i'\times k'\) matrices \(y\).  \emph{Status}: classical; see \cite{Hua}.  For one row, Problem~\ref{pr:hua-row} gives the exact integral for every \(k'\ge1\); its constant \(\pi^{k'}/k'!\) is at most \(1\) for \(k'\ge7\) (Problem~\ref{m2-pr:hua}).

\section{The size of a term}\label{sec:size}

Throughout, \(F\) and \(E\) are the fields of an admissible filling, with \(E\) containing a primitive \(5040\)-th root of unity; the data are such a filling.  For a term \(\gamma\) and classes \(c,c'\), write \(\tilde\gamma:=h_c^{-1}\gamma_fh_{c'}\); \(\gamma\) lies over the \emph{block} \((c,s,c')\) if \(\tilde\gamma\) is integral with reduction \(s\).

\begin{definition}\label{def:stable}
A set \(\Sf\) of terms is \emph{stable} if \(u^{-1}\gamma u'\in\Sf\iff\gamma\in\Sf\) for all \(u,u'\in\Uc(F)\).  An element \(s\in\bar G\) \emph{agrees} if the character of \(\Uc(\A_{F,f})\) of the definition of \(\Gamma^*\) in Chapter~\ref{m2-chap} takes the same value at any two integral \(u,u'\) in it whose reductions satisfy \(\bar us=s\bar u'\).
\end{definition}

\begin{definition}\label{def:cell}
A term \(\gamma\) has \emph{cell rank} \(j\in\{0,\dots,\varrho+1\}\) if (i) for some \(\zeta\in E\) with \(\zeta\bar\zeta=1\) and every index \(a\) with \(j<a\le\varrho+1\), some vector \(\sum_{b\le a}x_bv_b=\gamma\sum_{b\le a}y_bv_b\) has \(x_a=\zeta\), \(y_a=1\); (ii) \(\langle\gamma\rangle_j\ne0\); (iii) \(j=\varrho+1\) if and only if \(\disc(\gamma)\ne0\).  This rank, when it exists, is unique \cite[\texttt{BorderedDeterminant.lean}, l.~550]{GL}.  Write it as \(\cri(\gamma)\).
\end{definition}

\begin{definition}\label{def:near}
Let \(\Uc^0_c\) be the group of the \(u\in\Uc(F)\) that fix every frame vector and have \(h_c^{-1}u_fh_c\in K_M\).  The set \(\Rep(\Sf;c,s,c')\) of \emph{near representatives} of \(\Sf\) over \((c,s,c')\) consists of the \(\gamma\in\Sf\) over \((c,s,c')\) such that \(s\) agrees, \(\gamma\) has a cell rank, \(\gamma\) is the chosen representative of its double coset \(\Uc^0_c\gamma\,\Uc^0_{c'}\) \cite[\texttt{BorderedDeterminant.lean}, l.~455]{GL}, and \(\gamma\) is \emph{near}: some archimedean point \(y\) has \(y\) and \(\gamma_\infty^{-1}y\) boxed.  Use the domain \(D_0\) of \cite[\texttt{Box.lean}, ll.~145, 200]{GL}, a measurable set of representatives of the tuples \((\mathsf w_\tau)_\tau\) of complex numbers modulo the diagonal image of \(M\Oc_E\), with \(\lvert\mathsf w_\tau\rvert\le b^{26}\) on \(D_0\); a point \(y\) is \emph{boxed} if its Bessel length (Theorem~\ref{m1-thm:bessel-system}(d)) is at most \(b^{6(\varrho+1)+2}\) at every infinite place and each entry above the diagonal of its Gram--Schmidt matrix in the rational frame, as a tuple over the infinite places, lies in \(D_0\) \cite[\texttt{Setup.lean}, ll.~398--423]{GL}.  A \emph{pair of near representatives} over \(c,c'\) is a pair \((\gamma,\gamma')\) in \(\Rep(U(Y_0)(F);c,s,c')^2\) for some \(s\).
\end{definition}

\begin{definition}\label{def:calibre}
Let \(\gamma\) lie over a block and have cell rank \(\cri\).  Let \(\hat{\Oc}_E\) be the ring of integral finite adeles of \(E\), and let \(I(\gamma)\) be the index, in the \(\hat{\Oc}_E\)-module spanned by the rows \(a=1,\dots,\varrho\) of the matrix of products of the frame vectors \(v_a\) with the images of the frame under \(\tilde\gamma\), of the submodule spanned by its first \(\min(\varrho,\cri)\) rows, and \(I'(\gamma)\) the same index for \(\tilde\gamma^{-1}\); both are positive \cite[\texttt{LinkLattice.lean}, ll.~457, 464]{GL} and equal \(1\) if \(\cri\ge\varrho\), when all \(\varrho\) rows are taken.  The \emph{calibre} of \(\gamma\) is
\[
 \calb(\gamma):=M^{2f\min(\varrho,\cri)}\,\frac{\max\bigl(N^f\lvert\Nm_{E/\mathbb Q}\langle\gamma\rangle_{\cri}\rvert/N^{2f(\cri-\varrho)_+},\ \prod_\tau\lvert\tau(\disc\gamma)\rvert\bigr)}{I(\gamma)\,I'(\gamma)},
\]
where \(\tau\) runs over the \(f\) embeddings of the CM type.  For this formula see \cite[\texttt{Data.lean}, l.~74 and \texttt{BorderedDeterminant.lean}, l.~768]{GL}.  The frame product matrices used in \(I(\gamma)\) and \(I'(\gamma)\) are those of \(\gamma\) and \(\gamma^{-1}\) read in the finite adeles \cite[\texttt{Terms.lean}, l.~450]{GL}.  Thus \(\calb(\gamma)\) is independent of \(c,c'\).
\end{definition}

\begin{definition}\label{def:size}
The \emph{tail scale} is \(\ts:=b^{kf/4}\) and the \emph{pair scale} \(\pc:=\#\Rb\cdot n_{\mathrm{cl}}/N^{kf}\), with \(n_{\mathrm{cl}}\) the number of classes.  The \emph{cell moment} \(\Dcell(\gamma)\) is \(\Dmain\) if \(\cri(\gamma)=\varrho+1\) and \(\Did/\ts^{16}\) otherwise, and the \emph{size} of \(\gamma\) is \(\sz(\gamma):=\ts^2\bigl(\Dcell(\gamma)/(\pc\cdot\calb(\gamma)^k)\bigr)^{1/2}\).
\end{definition}

\begin{theorem}[size]\label{thm:size}
For a stable \(\Sf\subseteq\Gamma^*\), all classes \(c,c'\) and all \(s\in\bar G\),
\[
 \lVert\Amp^\flat(\Sf;c,s,c')\rVert\le\nu\sum_{\gamma\in\Rep(\Sf;c,s,c')}\sz(\gamma),
\]
a sum of nonnegative terms, possibly infinite.
\end{theorem}

\noindent\emph{Status}: read from the proof; not proved here.  For the Schur estimate on a double coset, see \cite[\texttt{SchurOrbit.lean}, l.~963]{GL}.

\section{Gain and count}\label{sec:gain}

\begin{problem}\label{pr:det-cs}
For \(x_1,\dots,x_i,y_1,\dots,y_i\in\C^{n'}\),
\[
 \lvert\det(\langle x_a,y_b\rangle)\rvert^2\le\det(\langle x_a,x_b\rangle)\det(\langle y_a,y_b\rangle).
\]
If all \(\langle x_a,y_b\rangle\) lie in \(M'\Z[\ii]\) for an integer \(M'\ge1\) and the determinant is not \(0\), the product of the two Gram determinants is at least \(M'^{2i}\).
\end{problem}

The gain at the level comes from a congruence: two terms over one block have the same reduction.

\begin{problem}\label{pr:congruence}
Let \(M'\ge1\) be an integer, \(\zeta_1\in\mathbb Q(\ii)\), \(\zeta_1\bar\zeta_1=1\), \(\zeta_1\ne1\), and let \(D'\ge1\) be an integer with \(D'(1-\zeta_1)\in M'\Z[\ii]\).  Then \(D'\ge M'/2\).  For \(\zeta_1=(-12+5\ii)/13\), \(M'=5\), \(D'=13\): \(13(1-\zeta_1)=5(5-\ii)\), \(\lvert5-\ii\rvert^2=26\), and \(\lvert1-\zeta_1\rvert\approx1.961\le2\).
\end{problem}

Over \(\Z[\ii]\) let \(h:=\diag(1,-5)\) and \(\gamma:=\left(\begin{smallmatrix}g_1&5g_2\\g_2&g_1\end{smallmatrix}\right)\) with \(g_1+g_2\sqrt5=(9+4\sqrt5)^3\).

\begin{problem}\label{pr:meet}
\(g_1+g_2\sqrt5=2889+1292\sqrt5\); \(\gamma\) is an isometry of \(h\) and \(\gamma\equiv-1\bmod17\); modulo each of the primes \(4\pm\ii\) above \(17\) it fixes the line of \(e_2\).  The determinant \(\det(e_2,\gamma e_2)=-6460\) lies in both primes; its norm \(6460^2=289\cdot144400\) is divisible by \(17\cdot17\), while \(6460\equiv102\bmod289\).
\end{problem}

In Problem~\ref{pr:meet} the two places above \(17\) at which the line meets its image contribute \(17\cdot17\) to the norm of the determinant, not to the determinant.

\begin{definition}\label{def:common}
Let \(\gamma,\gamma'\) lie over one block \((c,s,c')\), and let \(\mathrm{Lat}_c\) be the image under \(h_c\) of the lattice spanned by the lattice basis.  The \emph{common link} \(X_0(\gamma,\gamma')\subseteq X_0\) is the space of the \(x\in X_0\) for which some \(y\in X_0\) has \((x,y)\in\lk\gamma\cap\lk\gamma'\).  Lattices are taken in the finite adelic points \(V(\A_{E,f})\) of \(E\)-spaces \(V\), and products in the lattice coordinates of \(c\).  The \emph{term lattice} \(L_\gamma\) is the set of \(x\in\mathrm{Lat}_c\cap X_0(\A_{E,f})\) whose products with the frame of \(\tilde\gamma\) are \(M\) times the products of some vector of \(\mathrm{Lat}_c\cap(Ev_1+\dots+Ev_\varrho)(\A_{E,f})\) with it.  The \emph{common lattice} is \(L_{\gamma\gamma'}:=X_0(\gamma,\gamma')(\A_{E,f})\cap L_\gamma\cap L_{\gamma'}\).  If it is spanned over \(\hat{\Oc}_E\) by the \(\delta_0:=\dim X_0(\gamma,\gamma')\) columns of a matrix \(\mathsf Q=\mathsf x\mathsf u\), where \(\mathsf x\) consists of \(\delta_0\) linearly independent vectors of \(X_0(\gamma,\gamma')\) and \(\mathsf u\) is an invertible \(\delta_0\times\delta_0\) matrix of finite adeles, its \emph{discriminant} \(\dsc(\gamma,\gamma')\) is the square root of the inverse of the factor by which multiplication by \(\det(\mathsf Q^*\mathsf Q)\) scales Haar measure on the finite adeles, in orthonormal coordinates of \(X_0\); otherwise \(\dsc(\gamma,\gamma'):=0\).  The \emph{corank} is \(\crk(\gamma,\gamma'):=k-\dim X_0(\gamma,\gamma')\), and the \emph{link corank} is \(\lcr(\gamma,\gamma'):=k+1-\dim(\lk\gamma\cap\lk\gamma')\), which vanishes if and only if \(\lk\gamma=\lk\gamma'\).  Two terms are in the \emph{same link class} if both discriminants are nonzero and their links are equal.
\end{definition}

\begin{theorem}[gain]\label{thm:gain}
Let \(\gamma,\gamma'\) have cell ranks and lie over one block \((c,s,c')\) with \(s\) agreeing; put \(i:=\lcr(\gamma,\gamma')\), \(\ell:=\crk(\gamma,\gamma')\).
\begin{enumerate}
\item[(a)] For every integer \(t_0\) with \(0\le t_0\le\mt(s)\), let \(X:=t_0\) if \(i=0\) and \(\cri(\gamma)=\varrho+1\), and \(X:=1\) otherwise.  Then \((\dsc\cdot X)^2\le(2/M)^{2if}\calb(\gamma)\calb(\gamma')\).
\item[(b)] \(\ell\le i+(\cri(\gamma)-\varrho)_+\).
\item[(c)] If \(i=0\), then \(\disc(\gamma)=0\) if and only if \(\disc(\gamma')=0\).
\item[(d)] \(\dsc\ge1\); and if \((\gamma,\gamma')\) is a pair of near representatives, then \(\ell\le(\varrho+1)^3\) and \(\dsc\le b^{(\varrho+1)^3f}\).
\end{enumerate}
\end{theorem}

\begin{theorem}[count]\label{thm:count}
For every integer \(0\le\ell\le(\varrho+1)^3\), the sum of \(\dsc(\gamma,\gamma')^{-k}\) over all classes \(c,c'\) and the pairs of near representatives over \(c,c'\) of corank \(\ell\) with \(\dsc\le b^{(\varrho+1)^3f}\) is at most \(\ts^{\ell+2}\pc\).
\end{theorem}

\noindent\emph{Status}: both read from the proof; not proved here.  Problems~\ref{pr:det-cs}--\ref{pr:meet} show over \(\Z[\ii]\), as Problem~\ref{m2-pr:columns} does over \(\Z\), how a congruence modulo the level and a meeting place make a nonzero determinant large.  For a term \(\gamma\) with a cell rank over a block with \(s\) agreeing, \(\calb(\gamma)\ge1\).

\section{Pairs}\label{sec:pairs}

In this section \(\eta>0\), \(D\) and \(B\ge0\) are real numbers, \((\gamma,\gamma')\) is a pair of near representatives over \(c,c'\), \(i:=\lcr(\gamma,\gamma')\) and \(\ell:=\crk(\gamma,\gamma')\).

\begin{problem}\label{pr:pair-size}
Recall \(\dsc\ge1\) and \(\calb\ge1\).  If \((\dsc\cdot\eta)^2\le\calb(\gamma)\calb(\gamma')\) and \(\Dcell(\gamma),\Dcell(\gamma')\le D\), then \(\sz(\gamma)\sz(\gamma')\le\ts^4D/(\pc\,\eta^k\dsc^k)\).
\end{problem}

\begin{problem}\label{pr:rule}
For the data and all classes \(c,c'\), let \(\mathcal Q_{c,c'}\) be a set of pairs of near representatives over \(c,c'\), and suppose that every \((\gamma,\gamma')\in\mathcal Q_{c,c'}\) admits \(\eta,D\) as in Problem~\ref{pr:pair-size} with \(\ts^{\ell+6}D/\eta^k\le B\).  Then \(\sum_{c,c'}\sum_{(\gamma,\gamma')\in\mathcal Q_{c,c'}}\sz(\gamma)\sz(\gamma')\le\ts\cdot B\).
\end{problem}

\begin{problem}\label{pr:different-links}
For the data, \(\Dcell\le\Dmain\).  If \(i\ge1\) and \(\ell\le i+1\), then \(\ts^{\ell+6}\Dmain(M/2)^{-ifk}\le\ts^{-10}\Did\).  At \(i=1\), \(\ell=2\) the left side is at most \(b^{-35kf/9}\Did\), the right side \(b^{-10kf/4}\Did\).
\end{problem}

\begin{problem}\label{pr:zero-disc}
If \(\lk\gamma=\lk\gamma'\) and \(\gamma,\gamma'\) are not in the same link class, then \(\disc\gamma=\disc\gamma'=0\), \(\ell=0\), and \(\eta=1\), \(D=\Did/\ts^{16}\) give \(\ts^{\ell+6}D/\eta^k=\ts^{-10}\Did\).
\end{problem}

\begin{problem}\label{pr:same-class}
If \(\gamma,\gamma'\) are in the same link class over \((c,s,c')\) and \(1\le z_0\le\mt(s)\) is an integer, then \(\ell\le1\), and \(\eta=z_0\), \(D=\Dmain\) satisfy Problem~\ref{pr:pair-size} with \(\ts\cdot\ts^{\ell+6}\Dmain/z_0^k\le(b^{2f}/z_0)^k\Dmain\).
\end{problem}

Problems~\ref{pr:rule}--\ref{pr:same-class} with Theorem~\ref{thm:gain} prove the following.

\begin{theorem}[pairs]\label{thm:pairs}
For the data:
\begin{enumerate}
\item[(a)] the sum of \(\sz(\gamma)\sz(\gamma')\) over all classes \(c,c'\) and all pairs of near representatives over \(c,c'\) that are not in the same link class is at most \(b^{-kf}\Did\);
\item[(b)] for every integer \(z_0\ge1\), the same sum over the pairs in the same link class that lie over some \((c,s,c')\) with \(z_0\le\mt(s)\) is at most \((b^{2f}/z_0)^k\Dmain\).
\end{enumerate}
\end{theorem}

\section{The bounds of Chapter~\texorpdfstring{\ref{m2-chap}}{M2}}\label{sec:back}

Partition \(\Gamma^*\) into \emph{pieces} by labeling each term of nonzero discriminant by its link and giving all terms of discriminant \(0\) one further label.  We use the stability of each piece before restricting to a block \cite[\texttt{TermStructure.lean}, l.~843 and \texttt{Sum.lean}, l.~512]{GL}.  In each block the amplitude sum is the sum of the amplitude sums of the pieces \cite[\texttt{Sum.lean}, l.~498 and \texttt{Size.lean}, l.~709]{GL}.

\begin{problem}\label{pr:toy-pieces}
Four amplitudes \(2,-1,3,1\) with labels \(\Lambda,\Lambda,\Lambda',\varnothing\), where \(\Lambda\ne\Lambda'\) are links and \(\varnothing\) marks discriminant \(0\): the square of their sum is \(25\); the squares of the sums of the pieces \(\Lambda\), \(\Lambda'\) add to \(1+9=10\); the products \(\lvert a\rvert\lvert a'\rvert\) over the ordered pairs \((a,a')\) of amplitudes with different labels, or both labelled \(\varnothing\) (the amplitude \(1\) with itself), add to \(31\); and \(25\le10+31\).
\end{problem}

\begin{problem}\label{pr:pieces}
For all \(c,s,c'\), \(\lVert\Amp^\flat(\Gamma^*;c,s,c')\rVert^2\) is at most \(\sum_\Lambda\lVert\Amp^\flat(\Gamma^*\cap\{\disc\ne0,\ \lk=\Lambda\};c,s,c')\rVert^2\) plus \(\nu^2\) times the sum of \(\sz(\gamma)\sz(\gamma')\) over the pairs in \(\Rep(\Gamma^*;c,s,c')^2\) that are not in the same link class.
\end{problem}

\begin{problem}\label{pr:cut-moment}
For the data, let \(\Bc\subseteq\bar G\) and \(C_1\ge0\).  If the link second moment of \(\{\gamma:\disc(\gamma)\ne0\}\) over \(\Bc\) is at most \(C_1\nu^2\), then the second moment of \(\Gamma^*\) over \(\Bc\) formed with the cut function is at most \(C_1+b^{-kf}\Did\).
\end{problem}

\begin{problem}\label{pr:link-bound}
For the data, let \(z_0\ge1\) be an integer and \(\Bc\subseteq\bar G\) with \(z_0\le\mt(s)\) for every \(s\in\Bc\).  Then the link second moment of \(\{\gamma:\disc(\gamma)\ne0\}\) over \(\Bc\) is at most \((b^{2f}/z_0)^k\Dmain\,\nu^2\), which is Theorem~\ref{m2-thm:link-bound} for the data.
\end{problem}

\begin{theorem}[the tail]\label{thm:tail}
For every set \(\Sf\) of terms and every \(\Bc\subseteq\bar G\), the second moment over \(\Bc\) of the differences \(\Amp(\Sf;c,s,c')-\Amp^\flat(\Sf;c,s,c')\) is at most \(b^{-40}\Did\).
\end{theorem}

\noindent\emph{Status}: read from the proof; not proved here.

\begin{problem}\label{pr:cut}
For the data, let \(\Bc\) avoid the scalar cosets and \(C_0\ge0\), and let the link second moment of \(\{\gamma:\disc(\gamma)\ne0\}\) over \(\Bc\) be at most \(C_0\nu^2\).  Then the second moment of \(U(Y_0)(F)\smallsetminus R(F)\) over \(\Bc\) has square root at most \(\sqrt{C_0+b^{-kf}\Did}+b^{-20}\sqrt{\Did}\le\sqrt{C_0}+b^{-15}\sqrt{\Did}\).  The bound by \(\sqrt{C_0}+b^{-15}\sqrt{\Did}\) is Theorem~\ref{m2-thm:cut} for the data.
\end{problem}

\paragraph{What has been proved.}  Theorem~\ref{thm:pairs}, and from it the two bounds of Chapter~\ref{m2-chap}, Theorems~\ref{m2-thm:link-bound} and~\ref{m2-thm:cut} (Problems~\ref{pr:link-bound} and~\ref{pr:cut}), for the data.  Taken as given: Theorems~\ref{thm:size}, \ref{thm:gain}, \ref{thm:count} and~\ref{thm:tail}, read from the proof; Hua's integral (\S\ref{sec:one-place}); Theorem~\ref{m2-thm:survivors}; the two inequalities between \(\Did\) and \(\Dmain\); and stability and summation by pieces (\S\ref{sec:back}).

\begin{leanremark}
\raggedright In \cite{GL}, the following declarations are in the namespace \texttt{New.Filling}.
For Theorem~\ref{thm:size} see \tl{enorm\_amplitudeSum\_le\_tsum\_termSize}.
Theorem~\ref{thm:gain} specializes \tl{pairDiscriminant\_sq\_le} (a), \tl{le\_finrank\_commonLink\_add} (b), \tl{termDisc\_eq\_zero\_iff\_of\_link\_eq} (c), and \tl{one\_le\_pairDiscriminant} and \tl{IsNearPairOverAgreeingBlock.}\allowbreak\tl{corank\_le\_and\_pairDiscriminant\_le} (d).
For Theorems~\ref{thm:count} and~\ref{thm:tail} see \tl{sum\_pairDiscriminant\_inv\_le} and \tl{secondMoment\_sub\_cut\_le}, respectively.
Theorem~\ref{thm:pairs}(a), (b) specializes \tl{sum\_termSize\_mul\_not\_inSameLinkClass\_le} and \tl{sum\_termSize\_mul\_inSameLinkClass\_le} to the data.
\end{leanremark}

\section*{Hints}

\subsection*{\S\ref{sec:one-place}}
\begin{description}
\item[\ref{pr:schur}] Cauchy--Schwarz for \(x_j/\sqrt{v_j}\) and \(y_j\sqrt{v_j}\).
\item[\ref{pr:mean}] Unfold the sum over \(L\) against the integral over a fundamental domain; then polar coordinates (Problem~\ref{m2-pr:hua}).  \(\Z[\ii]\) has covolume \(1\) in \(\C\).
\item[\ref{pr:covolume}] The square of the covolume is a Gram determinant.
\item[\ref{pr:hua-row}] \(\det(A+v^*v)=\det A\,(1+vA^{-1}v^*)\); substitute \(v=\alpha y+y_0\), then \(v=wA^{1/2}\).
\end{description}

\subsection*{\S\ref{sec:gain}}
\begin{description}
\item[\ref{pr:det-cs}] The Gram matrix of \(x_1,\dots,x_i,y_1,\dots,y_i\) is positive semidefinite; or project the \(y_b\) to the span of the \(x_a\).
\item[\ref{pr:congruence}] A nonzero element of \(M'\Z[\ii]\) has absolute value at least \(M'\).
\item[\ref{pr:meet}] Factor \(6460\) and \(17\) in \(\Z[\ii]\).
\end{description}

\subsection*{\S\ref{sec:pairs}}
\begin{description}
\item[\ref{pr:pair-size}] Multiply the squares of the two sizes.  For \(\calb\ge1\), use Theorem~\ref{thm:gain}(a), (d) with \(\gamma'=\gamma\) and \(t_0=1\).
\item[\ref{pr:rule}] Sort the pairs by corank.  For the data \((\varrho+1)^3+1\le\ts\).
\item[\ref{pr:different-links}] Reserve one factor \((2/M)^{kf}\) for \(\Dmain\).  How large can \(\ts(2/M)^{kf}\) be?
\item[\ref{pr:zero-disc}] Theorem~\ref{thm:gain}(a) with \(t_0=1\), (b), (c), and Definition~\ref{def:cell}(iii).
\item[\ref{pr:same-class}] Theorem~\ref{thm:gain}(a), (b) with \(t_0=z_0\); \(\ts^4=b^{kf}\).
\end{description}

\subsection*{\S\ref{sec:back}}
\begin{description}
\item[\ref{pr:pieces}] Apply Theorem~\ref{thm:size} to the pieces.  Which squares must be kept in the expansion?
\item[\ref{pr:cut-moment}] Problem~\ref{pr:pieces}; a term over \(c,c'\) determines its reduction \(s\).
\item[\ref{pr:link-bound}] Theorem~\ref{thm:size} for one link, then Theorem~\ref{thm:pairs}(b).
\item[\ref{pr:cut}] Theorem~\ref{m2-thm:survivors}, Minkowski's inequality for the array of blocks, Theorem~\ref{thm:tail} and Problem~\ref{pr:cut-moment}.
\end{description}

\begin{thebibliography}{GL}
\bibitem[GL]{GL} L.~Sernau, \emph{GRH for Dirichlet \(L\)-functions, in Lean~4}, a shorter rewriting of the Lean proof, text by Claude, \url{https://github.com/lsernau/grh-lean}, commit \texttt{72fead1}, readable edition \texttt{grh/}.
\bibitem[Hua]{Hua} L.-K.~Hua, \emph{Harmonic Analysis of Functions of Several Complex Variables in the Classical Domains}, Transl. Math. Monographs \textbf{6}, Amer. Math. Soc., Providence, 1963.
\end{thebibliography}

\end{document}
